\documentclass[pre,twocolumn,superscriptaddress,nofootinbib]{revtex4-2}
\usepackage{graphicx}
\usepackage{amsmath}
\usepackage{amssymb}
\usepackage{booktabs}
\usepackage[hidelinks]{hyperref}

\begin{document}

\title{An entropy action across wirings: Bianconi's discrete entropic action evaluated on networks that rewire}

\author{Emily Smith}
\email{emilysmithcreate@gmail.com}
\affiliation{Independent researcher, Charlottesville, Virginia, USA}

\date{DRAFT, 10 October 2026. Series paper 10. Not reviewed by a physicist.}

\begin{abstract}
Bianconi's information-theoretic action for network geometry is defined on a cell complex with fixed wiring, and
its author names the dynamics of the wiring as open. We evaluate that action, unchanged, on networks whose wiring
does change: four-regular bipartite graphs from a family of square-counting energies next to combinatorial quantum
gravity, with every square taken as a 2-cell. (i) The construction passes the identities the source states.
(ii) In the action's vacuum the metric is the same on every cell, so the action is one number per cell and, across
wirings with the same number of nodes, a function of the number of squares alone; it orders wirings exactly as the
square-counting energy without its local term does. (iii) With the topology-dependent induced metric
$\mathbf{I}+c_0\boldsymbol{\mathcal{L}}$, evaluated at the identity metric, the action on 445 saved wirings of 64
nodes falls with every square and rises with every surplus square, as that energy does, with an effective local
coefficient $\lambda_{\rm eff}=0.001$ to 0.55 for $c_0=0.01$ to 100. That is below 1, the range in which curled
arrangements lie lower than the flat torus. The metric that solves the action's own equation for $c_0>0$ is not
computed, and whether a comparison across wirings is the intended one is put as a question to its author. Nothing
here is a statement about gravity.
\end{abstract}

\maketitle

\section{Introduction}

Ref.~\cite{Bia24} proposes an action for a network geometry coupled to matter and gauge fields: the quantum
relative entropy between a metric on the cells of a complex and the metric that matter induces on them. It is the
discrete form of ``gravity from entropy'' \cite{Bia25}. The complex is held fixed. The action ``depend[s] also on the
topology of the higher-order networks,'' and the paper leaves ``the possible implied dynamics of the network topology
to future works'' (its Sec.~3.1; its conclusions end on the same point).

The companion paper \cite{paper1} studies networks that do rewire, under an energy that counts squares. Its
arrangements can be handed to the action of Ref.~\cite{Bia24} as they are. We ask one narrow question: read in the
most direct way, what does that action say about arrangements with the same nodes and different wiring?

Three kinds of statement appear and are kept apart. \emph{Sourced}: taken from Ref.~\cite{Bia24}, with its equation
numbers. \emph{Exact}: linear algebra on named graphs, reproducible from the repository. \emph{Ours}: a reading or a
fit, unreviewed. The reading of the action across wirings is ours throughout.

\section{Setting}

\emph{The networks.} Four-regular bipartite graphs on $N$ vertices in which no two vertices share more than two
neighbors. With $S$ the number of 4-cycles (squares), $S_e$ the squares on link $e$ and $X=\sum_e(S_e-2)_+$ the
surplus squares, the energy is
\begin{equation}
H = 16(N-S) + 4\lambda X .
\label{eq:H}
\end{equation}
Only $\lambda=1$ is combinatorial quantum gravity (CQG) \cite{KTB19,T25}; other values are a family next to it. Three
arrangements of the same vertices are the rungs of a ladder \cite{paper1}: the flat torus ($S=N$, $X=0$), the curled
torus, four steps around in one direction ($S=5N/4$, $X=N$), and the 4-cube, curled both ways ($S=3N/2$, $X=2N$). At
$\lambda=1$ the three are degenerate; below 1 the curled ones lie lower; above 1 the flat torus is lowest.

\emph{The complex (ours).} Vertices, the $2N$ links, and every square as a 2-cell: $\mathcal{N}=3N+S$ cells. The
square count used here agrees with the simulation kernel's and with $(\mathrm{Tr}A^4-28N)/8$.

\emph{The action (sourced).} With boundary matrices $\mathbf{B}_{[1]},\mathbf{B}_{[2]}$ and a block-diagonal metric
$\boldsymbol{\mathcal{G}}$ on the cells, Ref.~\cite{Bia24} defines the weighted exterior derivative
$\mathbf{d}=\boldsymbol{\mathcal{G}}^{-1/2}\mathbf{B}^{\dagger}\boldsymbol{\mathcal{G}}^{1/2}$, the Dirac operator
$\mathbf{D}=\mathbf{d}+\mathbf{d}^{\dagger}$ (its Eq.~31), whose square is the Gauss--Bonnet Laplacian
$\boldsymbol{\mathcal{L}}=\mathbf{L}_{[0]}\oplus\mathbf{L}_{[1]}\oplus\mathbf{L}_{[2]}$ (Eqs.~35--40), and
\begin{equation}
\mathcal{S}_{+}=\sigma\,\mathrm{Tr}\ln\boldsymbol{\mathcal{G}}
+\mathrm{Tr}\,\boldsymbol{\mathcal{G}}\bigl(\ln\boldsymbol{\mathcal{G}}-\ln\mathbf{G}\bigr)
-\mathrm{Tr}\,\boldsymbol{\mathcal{G}}
\label{eq:S}
\end{equation}
(Eq.~41), with $\mathbf{G}$ the induced metric. With no matter and no gauge field,
$\mathbf{G}=\mathbf{I}+c_0\boldsymbol{\mathcal{L}}$ (Eq.~52), which the source notes is ``also dependent on the
topology.'' Its vacuum is $\mathbf{G}=\mathbf{I}$, where the metric obeys
$-\boldsymbol{\mathcal{G}}\ln\boldsymbol{\mathcal{G}}=\sigma\mathbf{I}$ (Eq.~65).

\section{The construction, checked}

Ref.~\cite{Bia24} reports no number to reproduce, so what is replicated is the construction. On every graph used
here (exact): $\mathbf{B}_{[1]}\mathbf{B}_{[2]}=0$; at the identity metric $\mathbf{D}^2$ equals
$\boldsymbol{\mathcal{L}}$ block by block; $\mathbf{D}$ anticommutes with $\gamma_0$ ($+1$ on vertices and squares,
$-1$ on links; Eq.~34); and the numbers of zero modes of $\mathbf{L}_{[0]},\mathbf{L}_{[1]},\mathbf{L}_{[2]}$ are
those of the complex: $(1,2,1)$ for the flat torus. For the curled torus they are $(1,1,N/4)$: each ring of four is
itself a square, a filled disk, so one independent loop survives and neighboring disks enclose $N/4$ cavities. A
4-cube gives $(1,0,7)$.

\section{The vacuum counts cells}

\emph{Exact, given the reading.} With $\mathbf{G}=\mathbf{I}$ and the same root $g$ of $-g\ln g=\sigma$ on every
cell, Eq.~(\ref{eq:S}) is
\begin{equation}
\mathcal{S}_{+}=\mathcal{N}\,s(\sigma),\qquad s(\sigma)=\sigma\ln g-\sigma-g .
\end{equation}
The wiring enters only through the number of cells. Since $\mathcal{N}=3N+S$ and $H_{\lambda=0}=16(N-S)$,
\begin{equation}
\mathcal{S}_{+}=4N\,s-\frac{s}{16}\,H_{\lambda=0},
\end{equation}
and $s<0$ (for example $g=0.894194$, $s=-1.005377$ at $\sigma=0.1$). At fixed $N$ the vacuum action is therefore an
increasing function of the square-counting energy with no local term: the two order every wiring alike. That energy
is the one under which the network is reported to break into isolated hypercubes \cite{KTB19,T25}. Two cautions. For
$\sigma<1/e$ Eq.~(65) has two roots, and a metric may take either on each eigen-direction; we take the root that
tends to 1 as $\sigma\to0$ everywhere. And ``more negative'' is not ``favored'' until a weight over wirings is
given, which the source does not give.

\section{At the identity metric}

For $c_0>0$ the induced metric carries the spectra. At
$\boldsymbol{\mathcal{G}}=\mathbf{I}$, where $\mathbf{d}=\mathbf{B}^{\dagger}$,
\begin{equation}
\mathcal{S}_{+}=-\sum_i\ln(1+c_0\mu_i)-\mathcal{N},
\label{eq:id}
\end{equation}
with $\mu_i$ the eigenvalues of the three Hodge Laplacians. This is the action at one uniform metric, not at a
solution of its equations (ours).

\begin{table}[t]
\caption{Four arrangements of 64 vertices. Zero modes are those of
$\mathbf{L}_{[0]},\mathbf{L}_{[1]},\mathbf{L}_{[2]}$. Vacuum: $\sigma=0.1$. Exact.}
\label{tab:ladder}
\begin{ruledtabular}
\begin{tabular}{lccccc}
 & $S$ & zero modes & vacuum & \multicolumn{2}{c}{Eq.~(\ref{eq:id})} \\
 & & & $\mathcal{S}_+$ & $c_0=0.1$ & $c_0=1$ \\
\hline
flat, $8\times8$ & 64 & 1, 2, 1 & $-257.38$ & $-339.46$ & $-642.03$ \\
curled, $16\times4$ & 80 & 1, 1, 16 & $-273.46$ & $-364.76$ & $-688.13$ \\
curled, re-glued & 80 & 1, 1, 16 & $-273.46$ & $-364.76$ & $-688.13$ \\
four 4-cubes & 96 & 4, 0, 28 & $-289.55$ & $-390.07$ & $-734.78$ \\
\end{tabular}
\end{ruledtabular}
\end{table}

\emph{The ladder (exact).} Table~\ref{tab:ladder}: the more squares, the more negative the action. The order,
4-cubes below curled torus below flat torus, holds at every $c_0$ tried from 0.01 to 1000.

\emph{A twin (exact).} Re-gluing two neighboring rings of the curled torus at opposite points gives a graph that is
not isomorphic to it and has a different spectrum (eigenvalues differ by up to 0.39), with the same cells. The two
agree in the first six moments of each Laplacian and in Eq.~(\ref{eq:id}) to nine figures: they differ only through
closed walks that go all the way round the long direction. At this size the action at the identity metric does not
see the global twist.

\emph{Why a penalty appears (ours).} For small $c_0$,
$\mathcal{S}_{+}=-\mathcal{N}-c_0\mathrm{Tr}\boldsymbol{\mathcal{L}}
+\tfrac12c_0^2\mathrm{Tr}\boldsymbol{\mathcal{L}}^2-\dots$, and
$\mathrm{Tr}\boldsymbol{\mathcal{L}}=8N+8S$: to first order the action counts squares and nothing else. At second
order pairs of squares that share a link enter with the opposite sign. Crowded links are penalized without that
being put in, and the surplus term of Eq.~(\ref{eq:H}) is one measure of crowding.

\begin{table}[t]
\caption{Fit of Eq.~(\ref{eq:id}) to $a+bS+cX$ over 445 wirings of 64 vertices: end states saved by the registered
runs of Ref.~\cite{paper1} (at most 25 of each kind $(S,X)$; 42 kinds, $61\le S\le96$) and the three rungs.
$\lambda_{\rm eff}=-4c/b$. Exact for each wiring; the fit is exploratory.}
\label{tab:fit}
\begin{ruledtabular}
\begin{tabular}{lcccc}
$c_0$ & $b$ & $c$ & $\lambda_{\rm eff}$ & rms left over \\
\hline
0.01 & $-1.078$ & $0.0002$ & 0.001 & 0.000 \\
0.1 & $-1.645$ & $0.015$ & 0.037 & 0.014 \\
1 & $-3.980$ & $0.250$ & 0.25 & 0.36 \\
10 & $-7.863$ & $0.891$ & 0.45 & 2.4 \\
100 & $-11.91$ & $1.62$ & 0.55 & 6.2 \\
\end{tabular}
\end{ruledtabular}
\end{table}

\emph{Many wirings (exact per wiring; the fit is ours).} Table~\ref{tab:fit}: every square lowers the action and
every surplus square raises it. Since $H=16(N-S)+4\lambda X$, the action orders these wirings as $H$ would with
$\lambda_{\rm eff}=-4c/b$, which grows with $c_0$ and stays below 1 (0.034, 0.20 and 0.28 at $c_0=0.1$, 1 and 10 if
the curled torus and the 4-cubes are left out). The linear form is a good description at small $c_0$ and a poorer
one as $c_0$ grows.

\section{Beside combinatorial quantum gravity}

In the family of Eq.~(\ref{eq:H}), $\lambda<1$ is where curling pays: over-full arrangements lie below the flat
torus. At $\lambda=1$ they tie, and above it the flat torus is lowest. Read at the identity metric, the action of
Ref.~\cite{Bia24} places these networks on the first side, for every $c_0$ tried. It contains the ingredient that
saturates the reward for squares, and too weakly to make flat space the floor.

What this does not show. That the comparison across wirings at fixed $N$ is the one the theory intends. Anything at
a metric that solves the equation of motion
$\sigma\boldsymbol{\mathcal{G}}^{-1}+\ln\boldsymbol{\mathcal{G}}=\boldsymbol{\mathcal{T}}$ (Eq.~57) for $c_0>0$,
where the metric will differ from cell to cell and from wiring to wiring. Anything with matter or gauge fields.
Anything about the continuum theory \cite{Bia25}. Anything about gravity.

\section{The question}

Is $\mathcal{S}_+$ compared across wirings at a fixed number of nodes the comparison meant by dynamics of the
topology? If it is, must the metric be solved for each wiring first, or is something else meant to keep the action
from rewarding more cells? The natural next steps follow from the answer: the metric equation solved on the ladder,
where symmetry leaves a few unknowns per rung; the same on a flat sheet holding one of the leftover knots of the
companion work; and six links per vertex.

\begin{acknowledgments}
Code, analysis and text were drafted with Claude (Anthropic) in conversation with the author, who directed and
checked the work. The passages of Ref.~\cite{Bia24} used here were read by an AI agent from the arXiv HTML and are to
be checked against the published version. No physicist has reviewed this. Everything is in the repository,
\url{https://github.com/EmilySmithCreate/RecreationalPhysics}: ASSUMPTIONS O111,
\texttt{scripts/exact\_entropy\_action.py}, \texttt{tests/test\_exact\_entropy\_action.py}.
\end{acknowledgments}

\begin{thebibliography}{9}
\bibitem{Bia24} G. Bianconi, ``Quantum entropy couples matter with geometry,'' J. Phys. A \textbf{57}, 365002 (2024),
arXiv:2404.08556.
\bibitem{Bia25} G. Bianconi, ``Gravity from entropy,'' Phys. Rev. D \textbf{111}, 066001 (2025), arXiv:2408.14391.
\bibitem{paper1} E. Smith, ``The curled torus burps: activated escape and conversion in a graph model of emergent
geometry,'' manuscript (2026), in the repository above.
\bibitem{KTB19} C. Kelly, C. A. Trugenberger and F. Biancalana, ``Self-assembly of geometric space from random
graphs,'' Class. Quantum Grav. \textbf{36}, 125012 (2019), arXiv:1901.09870.
\bibitem{T25} C. A. Trugenberger, ``Networks as the fundamental constituents of the universe,'' J. Phys. Complex.
\textbf{6}, 042001 (2025), arXiv:2512.17676.
\end{thebibliography}

\end{document}
